\section{Conclusions}
\label{sec:conclusions}

This paper asked whether any of the nearby stars that ALMA has observed shows
a narrowband, drifting radio carrier of the kind an artificial transmitter
would produce at millimetre and submillimetre wavelengths. Because the
observations were taken for other purposes, it also asked how deep and how
complete a search of the archive for such a carrier can be. To answer both
we searched the archival ALMA spectra of \NSysClassA{} star systems within
40\,pc. The fine-channel windows cover \CvUnionA\,GHz of sky frequency, in
\CovNIslands{} separate intervals between \SurvFreqLoA{} and
\SurvFreqHiA\,GHz. That union is \SbPriorRatio{} times the \SbPriorBandMHz\,MHz of the
wider of the two millimetre searches of nearby stars that preceded it, and
the coverage of the median system alone is still \SbPriorRatioSysMed{}
times larger. We removed the drift from each spectrum and compared it with
control positions in the same field, so that the background at the star
could be judged from the observation itself.

No transmitter was found. In \CsNCross{} windows the stellar statistic reaches
the threshold $T_\star\geq\StageNullTrig$. Of these, \MkNCoinc{} coincide with
catalogued molecular lines, \LnAttrSupp{} of them with further evidence of
line emission, and none was identified as interference by the available
diagnostics. The other \CeObs{} are consistent with chance: the
control positions beside each star predict \CeExp, with $P=\CeProbNull$. The extraction field and the line
mask could each have been chosen in two ways. Corrected for having been
picked out of those \CgNTrialWord{} combinations, the smallest of their
probabilities is \CgBonfProb, \CgVerdictClause. Both probabilities are conditional on the line mask as adopted, which was completed
after the crossings had been examined, and on the control positions being a
fair model of the background at the star.

None of the unattributed crossings became a recurrent feature, although
the molecular lines did, as they should. Wherever an unattributed
crossing was observed again in a separated session and the repeat could be
read, nothing was found at its own frequency in the star's frame. That shows
only that none of them returns in the way a line would. A drifting carrier
would by then have moved, and the archive can test one only by asking
whether a single trajectory joins two crossings, which can tell crossings
apart but can never confirm them. Recurrence is in any case necessary for an
artificial origin but not sufficient: a feature found again would still need
its line identity, interference and the behaviour of its frequency in the
star's frame to be checked. For the crossings themselves, therefore,
this is a non-detection, not an exclusion.

Because nothing was confirmed, what the search delivers is a set of
single-epoch upper limits. In the \XwNWin{} windows where an unattributed
crossing, or a weaker unattributed cell, is already present, a carrier
reaching only the threshold would not have stood out, so there the limit is
quoted at the statistic actually observed and is an upper limit as well. That moves the limits of \XwNSysWord{} systems, marked in
Table~\ref{tab:persys}, and \XwHeadlineClause. The limits apply to a
carrier narrower than one channel, lying outside the molecular-line mask,
radiating while ALMA observed and drifting by no more than
\DriftCeilLo--\DriftCeilHi\,nHz, the drift ceiling discussed in
\S\ref{sec:excludes}. To reach the threshold nine times in ten in a single
epoch, such a carrier needs an equivalent isotropic radiated power, or EIRP,
with a median, over systems, of $\EirpNinetySysMedA$\,W in each system's
most sensitive window. This is a single-epoch trigger completeness, not the
chance of confirming a signal, since a crossing would still have to be found
again. The deepest limit, $\EirpNinetySysLoA$\,W, is reached toward
$\epsilon$~Eridani, and \NSysEfifteen{} of the \NSysALimSum{} systems with a
limit that does not depend on an unchecked position reach
$\CovEirpRefTex$\,W.

It helps to set these numbers against something we have built ourselves.
The Arecibo planetary radar, one of the most powerful transmitters ever
operated on the Earth, had an EIRP of about $\AroEirp$\,W at
\AroFreqGHz\,GHz. The median limit found here is about \AroMedMult{} times
that, and even the deepest is \AroDeepMult{} times it, so
\AroNSysBelowWord{} system reaches an Arecibo-like transmitter, although
\AroNSysTen{} come within a factor of ten of it (counted on the unrounded
limits, \S\ref{sec:mmdifferent}). The same comparison can be
put as a distance. Because the received flux falls as the square of the
distance, an Arecibo-like transmitter placed at a star of this survey would
be detected only if the system were much closer to us than it is. For the
median system it would have to be within \AroDistMedPc\,pc, and even in the
most favourable case, \AroDistMaxStar, whose limit is deep for its distance,
within \AroDistMaxPc\,pc. With the atmospheric decorrelation divided out
these become \AroDistMedPcDec\,pc and, for \AroDistMaxStarDec, the most
favourable system on those values, \AroDistMaxPcDec\,pc, although \AroDistMaxStarDec's value rests on an
\AroDistMaxDecBasis{} decorrelation estimate. A \BenchDiam\,m dish radiating \BenchPowerMW\,MW is
a more natural transmitter at these frequencies. Tuned to the frequency of
each system's most sensitive window, between \BenchTunedFreqLo{} and
\BenchTunedFreqHi\,GHz, its EIRP lies above the single-epoch limit of
\BenchNSysTuned{} of the \NSysALimSum{} systems. It could only be seen, of
course, if its beam happened to point at us, as discussed in
\S\ref{sec:mmdifferent}.

These limits come from the single-epoch carrier search, and they are not the
deepest number in the paper. Toward Proxima Centauri the Class~B spectra,
stacked in the star's own frame, reach $\SzProxEirpCensus$\,W in
\ChanBModeMHz\,MHz channels (\S\ref{sec:stack}). That is a limit of a
different kind. It applies to an unresolved spectral excess that holds its
frequency in the stellar frame to within one of those wide channels over all
the epochs combined, it carries its own injection calibration, and it cannot
tell a carrier from a line. A transmitter on \PxPlanet{} would move across
many channels in its orbit, and the stack would cancel it unless the orbit
is seen nearly face-on, so for a drifting carrier the single-epoch limits
remain the ones that apply.

How complete the search can be is set more by the archive's coverage than by
its depth. A crossing can be put to the test only if its star is observed
again in a separated session, at the frequency the crossing predicts. Of the
\NSysClassA{} systems, \ZcNoRepNoSep{} were never observed in two separated
sessions at all, so that nothing found toward them could have been confirmed
at any power. For \ZcNoRepNoFreqWord{} more, separated sessions exist but
share no covered frequency, so the same-frequency recurrence test cannot be
applied there. A programme designed for the purpose, on the same telescope,
would use fine channels throughout, tune away from the bright molecular
lines and schedule its repeat observations deliberately, so that any
crossing it found could be looked for again at the frequency it predicts
(\S\ref{sec:future}).

What, then, can be said about transmitters? A continuous, unresolved carrier
at one of these stars, outside the line mask, drifting within the ceiling and
radiating above the limit of the system's most sensitive window would have
crossed the threshold nine times in ten in a single epoch. Such a
transmitter is therefore excluded at the single-epoch trigger level, toward
the systems and at the frequencies that the archive covers. A great deal
remains possible. A transmitter below these limits would not have been seen,
and neither would one that was switched off, or whose beam was not pointed
at the Earth, while ALMA was looking. Signals at frequencies between the
covered intervals or inside the line mask, drifting faster than the ceiling
or broader than a channel, are not constrained at all. Nor is anything said
about the nearby stars that ALMA has not observed, which are almost all of
those within 40\,pc.
